\documentclass[11pt,a4paper]{article}

% ========== パッケージ群 ==========
\usepackage{amsmath}
\usepackage{amssymb}
\usepackage{amsfonts}
\usepackage{graphicx}
\usepackage{cite}
\usepackage{hyperref}
\usepackage{url}
\usepackage{xcolor}
\usepackage{geometry}
\geometry{margin=1in}

% ========== 定理環境 ==========
\usepackage{amsthm}

\theoremstyle{definition}
\newtheorem{lemma}{補題}[section]
\newtheorem{definition}{定義}[section]
\newtheorem{proposition}{命題}[section]
\newtheorem{assumption}{仮定}[section]
\newtheorem{remark}{注記}[section]
\newtheorem{phenomenological}{現象論的原理}[section]
\newtheorem{clarification}{説明}[section]

% ========== XeLaTeX 用日本語対応（macOS） ==========
\usepackage{xeCJK}
\setCJKmainfont{Hiragino Sans}
\setCJKsansfont{Hiragino Sans Mono}
\setCJKmonofont{Monaco}

% ========== ハイパーリンク設定 ==========
\hypersetup{
  colorlinks=true,
  linkcolor=blue,
  citecolor=blue,
  urlcolor=blue,
  pdftitle={意味の非収束と共鳴条件},
  pdfauthor={加納 智之},
  pdfsubject={Resonanceverse 理論の基礎},
  pdfkeywords={共鳴, 誤配, 意味動力学, 非収束安定性},
}

% ========== バージョン管理 ==========
\newcommand{\versionnum}{1.1}
\newcommand{\versiondate}{2026/05/03}

% ========== タイトル情報 ==========
\title{意味の非収束と共鳴条件：\\
誤配に基づく Resonanceverse の基礎理論\\
\large \textit{Version \versionnum, \versiondate}}

\author{加納 智之$^{1,*}$}

\date{\today}

% ========== 本文開始 ==========
\begin{document}

\maketitle

% ========== 所属・連絡先（タイトル直下） ==========
\begin{center}
\small
$^{1}$ ZYX Corp 株式会社 人工叡智研究室（Artificial Sapience Laboratory）\\
$^{*}$ 責任著者（Corresponding author）\\
E-mail: \url{tomyuk@zyxcorp.jp}\\
ORCID: \url{https://orcid.org/0009-0004-8213-4631}
\end{center}

\vspace{0.5cm}

\begin{abstract}
本論文は Resonanceverse を提示する。これは意味を誤配（misdelivery）に条件付けられた非収束的動力学プロセスとして再定義する理論的枠組みである。ポスト構造主義、特に Jacques Derrida は意味の不安定性を開示したが、そうした洞察は動力学モデルと結合されない限り、概ね記述的なままである。

従って、意味は $\mathbb{R}^d$ 内の符号化された状態軌道としてモデル化される。ここで差異と規範が十分に定義されている。形式的には、$\Delta M_{ij}(t) := M_j(t) - M_i(t)$ が二つのエージェント $(i, j)$ 間の意味論的状態差を時刻 $t$ で表す。ここで $M_k(t) \in \mathbb{R}^d$ である。

誤配は $\Delta M$ と同一ではない。むしろそれは、相互作用論的な生成的層として機能し、通じてそうした差異が生成され、持続し、チャネルを通じて再形成される。それでも $\Delta M$ は、状態がいったん符号化されたら、定量的層でどのように不整合が進化するかをインデックスする。

関係は文脈的重みと差異を結合する時間積分としてモデル化される。共鳴には意味的、不確実性、および価値制約 $C_M(t), C_U(t), C_V(t)$ の同時満足が必要である。本論文は、意味的アイデンティティへの収束なしに、実用的安定性がいかにして現れるかを明確にする。

本アプローチは意味を固定的な実体ではなく、継続的な偏差の持続的でかつ構造的に制約された過程として再定義する。ポスト構造主義哲学と動力学系理論を架橋することにより、Resonanceverse は人工知能、通信理論、制度設計への応用のための形式的基礎を提供する。

包括的な数値シミュレーション検証により、全ての理論的予測が確認される。Phase 1～4 の実験は非収束安定性、制度的コンセンサス形成、共鳴条件の満足（タイムステップの 98.6\%）、および非収束公理の反証可能性を実証する。
\end{abstract}

\section*{キーワード}
共鳴、誤配、意味動力学、非収束安定性、マルチエージェント $\Delta M$ ネットワーク、コンセンサス、共鳴条件、ポスト構造主義、ジャック・デリダ、動力学系、制度的制約、ゲーデル類推、Resonanceverse

\section{導入}

意味とは何か。この問いは哲学、言語学、認知科学、情報科学において何度も提起されてきた。しかし多くの理論は、明示的であれ暗黙的であれ、意味を比較的に安定した構造として、あるいは最終的に同定可能な状態として前提している。

この前提はポスト構造主義思想によって根本的に再検討された。とりわけ Jacques Derrida は、意味が自己同一性を持たず、\textit{différance} として延期され、散在することを示し、それゆえその不確定性と非閉鎖性を開示した。東浩紀はさらにこの不確定性を「郵便的脱構築」として再解釈し、意味が誤配を通じて生成されるというモデルを提案した。

しかし、こうした議論は概ね記述的・批判的枠組みに留まっている。それらは、意味の時間的変換と相互作用的動力学を記述できる形式的理論にまで至っていない。言い換えれば、以下の問いは十分に形式化されていない。

\textit{意味はいかにして時間の中で変化し、持続し、崩壊するのか？}

本研究は意味動力学の枠組みを導入することでこの隙間に対処する。意味を時間的に変動する状態量として扱う。形式的モデルでは、各エージェント $i$ に対して意味状態 $M_i(t)$ が割り当てられ、その差異
\[
\Delta M_{ij}(t) := M_j(t) - M_i(t)
\]
は二つのエージェント間の意味論的状態差として扱われる。この $\Delta M$ は誤配そのものと同一ではない。むしろ、チャネル全体での偏差、変位、再配送といった誤配プロセスが $\Delta M$ を時間的に生成し、増幅し、あるいは減衰させる。誤配プロセスは実用的な生成的層として機能し、一方 $\Delta M$ はそのプロセスが動力学系にマップされたときの基本的な観測量として機能する。このポジションはエラーを単なるノイズとして扱うアプローチとは対照的である。

\subsection{中心的仮説}

本論文の中心的仮説は以下の通りである。

\textit{意味は正しく伝達されるから持続するのではなく、継続的に誤配されるがゆえに非収束的な動力学的構造として持続する。}

この視点から、意味は静的なオブジェクトではなく、関係的かつ時間的なプロセスとして定義される。意味は単一のエージェント内部に存在しない。むしろそれは複数のエージェント間の相互作用を通じて生成される。従って、意味変化はエージェント間の差異のネットワーク、すなわち $\Delta M$ ネットワークとして記述される。

これは従来のコンセンサス形成理論に仮定されたコンセンサスの概念の再検討を要求する。本論文ではコンセンサスを意味の同一性ではなく、エージェント間の $\Delta M$ の分布が有界で安定した構造を形成する状態として再定義する。この意味において、コンセンサスは偏差の消滅ではなく、偏差の構造的修正である。

さらに本論文は意味的安定性を収束ではなく、非収束下での有界性として再定義する。この枠組みでは、意味は原理的に収束せず、継続的に変動する。しかし、そうした変動が特定の構造的条件を満たすとき、システム全体は安定した振る舞いを示す。このような状態は\textit{共鳴}と呼ばれる。

\subsection{論文の貢献}

このフレームワークに基づいて、本論文の貢献は五つである。

\begin{enumerate}
\item 意味の非固定的かつ非収束的性質を公理系として形式化する。
\item 誤配プロセスを生成的層として位置付け、$\Delta M$ ネットワークの動力学を対応する動力学系の観測可能層として導入する。
\item コンセンサス形成をマルチエージェント $\Delta M$ ネットワークの形成として再定義する。
\item 共鳴条件を非収束下での安定性の条件として提示する。
\item 制度設計およびエージェント・アーキテクチャへの応用可能性を示唆し、包括的な数値検証により支持する。
\end{enumerate}

本研究はポスト構造主義の洞察を否定しない。むしろ、そうした洞察を保存しながら、意味理論を記述から動力学系理論へ移動させることを目指している。この再定義は意味を固定的なオブジェクトではなく、関係の内部に持続する動的現象として扱う。それゆえ、インテリジェンスのモデル化、人工知能システムの設計、および社会制度の理論化のための新たな基礎を提供している。

\section{意味動力学の公理系}

本セクションは意味を時間的に変動する状態量として扱うための公理系を定義する。本研究では、意味は固定的な表現や内在的属性ではなく、エージェント間の相互作用を通じて継続的に更新される動的状態である。従って意味は単一のポイントへの収束として記述されるのではなく、時間を通じて移動する状態として記述される。

時刻 $t$ にエージェント $i$ が保有する符号化された意味状態を $M_i(t) \in \mathbb{R}^d$ で表す。$\mathcal{M}$ を抽象的な意味状態空間とする。測定可能な状態は符号化マップ
\[
\phi : \mathcal{M} \to \mathbb{R}^d
\]
を通じて $\mathbb{R}^d$ で表現されていると仮定される。

記述の簡潔性のため、符号化された像を単に $M_i(t)$ で表記する。差異と規範は $\mathbb{R}^d$ の通常の線形代数的構造を用いて定義される。距離は有するが減法をもたない一般的な計量空間では $\Delta M_{ij} = M_j - M_i$ を記述しない。

二つのエージェント $i$ と $j$ 間の符号化された状態差は次のように定義される。
\[
\Delta M_{ij}(t) := M_j(t) - M_i(t) \in \mathbb{R}^d.
\]

形式的には、$\Delta M_{ij}(t)$ は二つのエージェントが保有する意味状態間の差異である。従って $\Delta M$ は誤配そのものと同定されるべきではない。むしろ誤配はチャネル、制度、規範の実用的側にあることもある。本論文では、$\Delta M$ を誤配プロセスを通じて現れ、持続し、変換される意味的差異として制約付けで定義する。この区別は理論の数学的および修辞的構造の枢軸となる。

本論文では $\Delta M$ を多次元ベクトルとして扱う。その大きさは従ってノルムで測定される。明示されない限り、ユークリッド規範が使用される。
\[
\|\Delta M_{ij}(t)\| := \left(\sum_{k=1}^{d} (\Delta M^{(k)}_{ij}(t))^2\right)^{1/2}.
\]

各成分は意味、価値、不確実性など異なる次元に対応することがあり得る。

従来の理論では、そのような量は誤解またはノイズとして扱われることがある。本研究では、それを意味変化の基本的単位として扱う。意味は恒等性ではなく差異によって駆動される。そのような差異は例外的ではない。それは通常の条件であり、意味生成の出発点である。

\subsection{核心的公理}

この前提に基づいて、最初の公理は意味の時間的性質に関わる。全てのエージェントについて、意味状態は時間上で不変ではない。
\[
\forall I \subset [0,\infty), \mu(I) > 0 \Rightarrow \exists t_1, t_2 \in I \text{ such that } M_i(t_1) \neq M_i(t_2).
\]

この公理は意味が完全に静的なままではないことを意味する。正の測度を持つ任意の区間においても、意味は少なくともわずかに変動する。完全な恒等状態は時間全体を通じて持続しない。

さらに、意味生成の過程は孤立的ではなく、他者との相互作用に依存する。これは次のように表現される。
\[
M_i(t) = M_i(0) + \int_0^t \sum_j f(\alpha_{ji}(\tau), \Delta M_{ji}(\tau)) d\tau,
\]
ここで関数 $f$ は文脈的親和性と意味的差異をエージェント $i$ の意味状態の変化ベクトルへマップする。
\[
f : [0, 1] \times \mathbb{R}^d \to \mathbb{R}^d.
\]

この表式はエージェント $i$ の意味状態がその初期状態に他のエージェントとの関係で生じる意味的差異の累積を加えたものから構成されることを述べている。それは意味が内部生成ではなく関係的プロセスとして形成されることを明確にする。言い換えれば、意味は他者からの偏差の歴史である。

\subsection{非収束性公理}

次に、次の公理は意味的差異の時間的極限に関わる。
\[
\limsup_{t \to \infty} \|\Delta M_{ij}(t)\| > 0.
\]

この非収束性公理は、エージェント間の意味的差異は時間上で消滅しないと述べている。どれほど相互作用が継続しても、原理的に完全な恒等性は達成されない。この点は意味の最終的統一を仮定するモデルとは区別される。

\begin{lemma}[スケール分離]
非収束性公理 $\limsup_{t \to \infty} \|\Delta M_{ij}(t)\| > 0$ は軌道がゼロの任意小近傍に入ることを禁止しない。それが否定するのは差異の完全な消滅、すなわち厳密な恒等性である。

形式的には、$\varepsilon_0 > 0$ と $T \geq 0$ が存在して
\[
\forall t \geq T: 0 < \|\Delta M_{ij}(t)\| < \varepsilon_0
\]
が同時に成立しうるが、一方で
\[
\limsup_{t \to \infty} \|\Delta M_{ij}(t)\| > 0
\]
が成り立つ。この二層構造——論理的層での「ゼロに向かわない」と実用的層での「$\varepsilon$ スケールで修正可能」——は\textit{スケール分離}と呼ばれる。
\end{lemma}

意味的差異の累積はエージェント間の関係を構成する。関係は以下のように定義される。
\[
R_{ij}(t) = \int_0^t g(\alpha_{ij}(\tau), \Delta M_{ij}(\tau)) d\tau,
\]
ただし $g : [0, 1] \times \mathbb{R}^d \to \mathbb{R}^r$ である。

関数 $g$ は文脈的親和性と意味的差異を関係状態へマップし、その値域は一般に $\mathbb{R}^r$ である。この定義の下では、関係は静的な属性ではなく、過去にどのような意味的差異が生じ、いかにして累積されたかの歴史である。たとえ二つのエージェントが現在の意味状態を共有しても、異なる関係的歴史によって区別されることがあり得る。

しかし、全ての意味的差異が同じように作用するわけではない。エージェント間の文脈的親和性は意味的変化の影響を調整する。$\alpha_{ij}(t)$ をこの文脈的親和性とし、$\alpha_{ij}(t) \in [0, 1]$ とする。

実効的意味変化は次のように定義される。
\[
\Delta M^{\text{eff}}_{ij}(t) = \alpha_{ij}(t) \Delta M_{ij}(t).
\]

この方程式は意味的差異が直接的に意味的更新に貢献するのではなく、関係と文脈に従って重み付けされることを示す。意味的変化は従ってただ差異の関数ではなく、関係によって仲介されたプロセスである。

意味状態の時間進化は次のように記述される。
\[
\frac{dM_i(t)}{dt} = \sum_j f(\alpha_{ji}(t), \Delta M_{ji}(t)).
\]

ここで $f$ は一般に非線形であり、歴史的依存を含むことがある。この定式化は意味が単純な加算的プロセスではなく、複雑な相互作用システムとして振る舞うことを示す。

マルチエージェントシステム全体の水準では、この相互作用はネットワークとして表現される。$G(t) = (V, E(t))$ を、エージェント集合 $V$ と相互作用集合 $E(t)$ から構成されたグラフとする。各辺は二つのエージェント間の意味的差異 $\Delta M_{ij}(t)$ に対応する。意味は個々のエージェント内に閉じ込められていない。むしろそれはネットワーク全体に分布している。

そのようなシステムにおいて重要な点は、意味的差異が消滅していなくても、システムが安定していることがあり得ることである。本論文は安定性を次の条件として定義する。
\[
\exists T \geq 0, \exists \varepsilon > 0 \forall t \geq T, \forall (i,j) \in E(t) : 0 < \|\Delta M_{ij}(t)\| < \varepsilon.
\]

この条件は意味的差異が存在し続ける一方、無限に発散する代わりに一定範囲内で有界であることを意味する。意味は収束なしに安定化する。この非収束安定性は古典的安定性概念に基づく収束とは異なった概念的枠組みを提供する。

最終的に、時間の経過と共に消滅しない残存する意味的差異が導入される。グラフ $G(t)$ と混同を避けるため、残存は $G_{ij}(T)$ と書かれ、
\[
G_{ij}(T) := \sup_{t \geq T} \|\Delta M_{ij}(t)\|
\]
である。

この量はインタラクションを通じて完全には解決されず、長期的振る舞いに影響し続ける差異を表現する。残存は意味が完全には統一されず、一方で過去の相互作用が現在に影響し続ける機構を提供するという事実の数学的表現である。

\subsection{説明的二エージェント例}

公理的枠組みを説明するため、次のパラメータを持つ $\mathbb{R}^2$ の単純な二エージェントシステムを考えよう。

\begin{itemize}
\item 初期状態：$M_1(0) = [0.0, 0.0]$、$M_2(0) = [1.0, 0.5]$
\item 初期差異：$\Delta M_{12}(0) = [1.0, 0.5]$、$\|\Delta M_{12}(0)\| \approx 1.12$
\item 相互作用関数：
\[
f(\alpha, \Delta M) = \begin{cases}
-0.5 \alpha \Delta M & \text{if } \|\Delta M\| < 0.1 \text{ (短距離反発)} \\
0.3 \alpha \Delta M & \text{if } \|\Delta M\| \geq 0.1 \text{ (長距離引力)}
\end{cases}
\]
\item 文脈的親和性：$\alpha_{12}(t) = 0.7 + 0.2\sin(0.1t)$（時間変動）
\end{itemize}

\textbf{期待される動力学：}
\begin{enumerate}
\item \textbf{早期フェーズ}（$t \approx 0$～20）：大きな意味的差異はエージェント間を長距離引力を通じて駆動する。$\|\Delta M_{12}\|$ は単調に $1.12$ から $\approx 0.3$ に減少する。

\item \textbf{中期フェーズ}（$t \approx 20$～60）：エージェントが近づくにつれて、短距離反発が開始される。振動的振る舞いが現れる。$\|\Delta M_{12}\|$ は $0.1$ 周辺で減衰振動を示す。

\item \textbf{後期フェーズ}（$t \approx 60$～100）：システムは有界非収束レジームに入る。$\|\Delta M_{12}\|$ は $[\varepsilon_0, \varepsilon_1] = [0.05, 0.15]$ 内で無期限に変動する。
\end{enumerate}

この例は数値的に補題1を検証する。意味的差異は非零のままで有界に保たれ、スケール分離が実践においていかに作用するかを実証している。

この公理系を通じて、意味は個々のエージェントの属性ではなく、エージェント間の差異と彼らの歴史から構成される動的構造として形式化される。本セクションで提示された枠組みは次のセクションでのマルチエージェント $\Delta M$ ネットワークにおけるコンセンサス形成の再定義へと発展する。

\section{マルチエージェント $\Delta M$ ネットワークとコンセンサス形成}

本セクションはセクション2で定義された意味動力学をマルチエージェントシステムに拡張し、コンセンサス形成を再定義する。従来のコンセンサス理論はしばしば暗黙的に、エージェント間の意味が一致する状態、すなわち $M_i = M_j$ を仮定している。しかし、セクション2で示されたように、意味は非収束的であり、意味的差異 $\Delta M_{ij}(t)$ は原理的に消滅しない。

この前提の下では、完全な恒等としてのコンセンサスは一般的に成立しない。

この問題に対処するため、本研究はコンセンサスを意味の恒等性ではなく、意味的差異の構造として再定義する。コンセンサスは、指定された条件下で $\Delta M$ ネットワーク間のエージェントが安定した分布を形成する状態である。

エージェント集合を $\mathcal{A} = \{a_1, a_2, \ldots, a_n\}$ とする。すなわち $\mathcal{A}$ は意味的相互作用ネットワークに参加するエージェントの有限集合を表す。

各エージェントは符号化された状態 $M_i(t) \in \mathbb{R}^d$ を持つ。意味的差異がその成分である構造行列 $X(t)$ を $X(t)$ で表す。マトリックスレンダリングの不安定性を避けるため、成分ごとに定義される。
\[
[X(t)]_{ij} := \Delta M_{ij}(t) \quad (i \neq j), \quad [X(t)]_{ii} := 0.
\]

この構造 $X(t)$ はマルチエージェントシステムの意味的差異ネットワークを表現する。ここで重要な点は意味が個々のエージェント内に限定されず、この差異のネットワークとして存在することである。意味は分散し、エージェント間の関係を通じてのみ定義される。

ネットワークの時間進化はセクション2で導入された動力学に従う。
\[
\frac{dM_i(t)}{dt} = \sum_j f(\alpha_{ji}(t), \Delta M_{ji}(t)).
\]

このインタラクションを通じて、意味的差異の分布は時間の中で変化するが、必ずしもゼロに収束しない。実際のマルチエージェントシステムでは、持続的な意味的差異は一般に規則ではなく例外である。

この視点に基づいて、コンセンサスは以下のように定義される。

\[
\text{コンセンサス}(t) \Leftrightarrow \forall (i,j) \in E(t) : \|\Delta M_{ij}(t)\| < \varepsilon \land -\delta < \frac{d}{dt}\|\Delta M_{ij}(t)\| < \delta.
\]

許容値 $\varepsilon$ はセクション2の非収束性公理 $\limsup_{t \to \infty} \|\Delta M_{ij}(t)\| > 0$ と矛盾しない。むしろそれは論理的層で差異の完全な消滅を除外しつつ、実用的許容帯を標示する。これはスケール分離として前に議論されたものである。

ここで $\varepsilon > 0$ と $\delta > 0$ は許容閾値である。後者の制約は一方向的条件 $\frac{d}{dt}\|\cdots\| < \delta$ ではなく、上記に示された二方向不等式である。それは規範の時間導関数が $-\delta$ と $\delta$ の間に留まることを要求し、それゆえ急速な増加と急速な減少の両方を抑制する。

この構造は分散の観点からも記述されることがある。平均エッジ規範を
\[
\mu_\Delta(t) = \frac{1}{|E(t)|} \sum_{(i,j) \in E(t)} \|\Delta M_{ij}(t)\|
\]
とする。

エッジ上の規範の分散は
\[
\text{Var}_E(\Delta M(t)) = \frac{1}{|E(t)|} \sum_{(i,j) \in E(t)} (\|\Delta M_{ij}(t)\| - \mu_\Delta(t))^2
\]
である。

この量はマルチエージェントシステムにおける意味的差異規範の広がりを表現する。分散が大きい場合、意味的差異は広く分散し、システムは断片化されていると解釈されることがある。分散が小さすぎる場合、意味的差異はほぼ不在であり、システムは均質化される。そのような状態は安定して見えるかもしれないが、外部変化に対して脆弱である。

従ってコンセンサスはゼロへ向かう分散ではなく、安定した中間範囲によって特徴付けられる。この状態は意味的差異が存在し続ける一方で、それらの変動が制御されることを意味する。

このように定義されたコンセンサスは従来の恒等性モデルと異なり、差異の存在を前提とする。コンセンサスは偏差の消滅ではなく、偏差の構造的修正である。

最後に、この枠組みにおける制度の役割に対処する必要がある。制度はエージェント間の相互作用を制約し、それゆえ $\Delta M$ の分布を制御する。制度は意味そのものを固定しない。それはエージェント間の差異の範囲と変動を制約する。この視点から、制度は次のように表現されることがある。
\[
\text{制度} = \text{$\Delta M_{ij}(t)$ への制約}
\]

この定義は制度を静的規範ではなく動的制御メカニズムとして扱う。

制度設計は意味的差異の分布がいかに制御されるべきであり、それがいかなる範囲内で安定化されるべきかの問題である。

従ってマルチエージェントシステムにおける意味はエージェント間の差異のネットワークとして記述され、コンセンサスはそのネットワークにおける安定した構造として現れる。本セクションで提示された枠組みは次のセクションで導入される共鳴条件を通じて非収束性下での安定性の、より精密な定義へ発展する。

\section{共鳴と非収束安定性}

本セクションはセクション2および3で定義された意味動力学とマルチエージェント $\Delta M$ ネットワークに基づいて、非収束下での安定性の条件として共鳴を導入する。従来の理論では、安定性はしばしば収束を通じて定義される。すなわち、意味状態が固定点に到達することによってである。本研究の非収束性前提の下では、そのような定義を適用することができない。従ってそれは意味が収束しないにもかかわらず安定した振る舞いを示す状態を定義することが必要である。

意味的差異 $\Delta M_{ij}(t)$ が時間上で消滅しない状態を考え、一方その大きさは一定範囲内で留まる。
\[
\Delta M_{ij}(t) \neq 0 \land \|\Delta M_{ij}(t)\| < \varepsilon.
\]

この条件は意味的差異が存在し続ける一方で無限に発散する代わりに制約されている状態を表現する。そのような状態では、エージェント間の意味は完全に一致することはないが、差異の範囲が安定しているため、システム全体は秩序立った振る舞いを示す。この安定性は収束ではなく有界性によって特徴付けられる。

しかし、この条件だけでは不十分である。意味的差異の単なる有界性は周期的振動または混沌的振る舞いを含むことがある。従ってその規範の時間導関数も実用的閾値によって制約される必要がある。
\[
-\delta < \frac{d}{dt}\|\Delta M_{ij}(t)\| < \delta.
\]

これはセクション3のコンセンサス定義に同じ $\delta$ である。それは規範の変化が実用的に受け入れ可能なほど遅く留まる状態を指定する。

これらの条件を満たす状態は\textit{非収束安定性}と呼ばれる。それは意味が固定点に収束することなく有界領域内で変動し続ける状態を指す。重要な点は、変動そのものが不安定を暗示しないこと。むしろ、安定性は変動が制約された範囲内で維持されるがゆえに、ちょうどそこで生じることがあり得ることである。

この非収束安定性をより精密に特徴付けるため、本論文は\textit{共鳴}の概念を導入する。記号 $C_M(t), C_U(t), C_V(t)$ は各時刻で真理値を割り当てる述語または制約を表現する。それらは意味状態ベクトル $M_i(t)$ の省略形ではない。一つの具体的なインスタンス化として、閾値 $\theta_M, \theta_U, \theta_V$ を用いて以下を定義する。
\begin{align}
C_M(t) &: \Leftrightarrow \text{Var}_E(\Delta M(t)) < \theta_M, \\
C_U(t) &: \Leftrightarrow D_U(t) < \theta_U, \\
C_V(t) &: \Leftrightarrow D_V(t) < \theta_V.
\end{align}

共鳴の概念そのものは論理積である。
\[
\text{共鳴}(t) \Leftrightarrow C_M(t) \land C_U(t) \land C_V(t).
\]

$C_M, C_U, C_V$ が上記の閾値定義によってインスタンス化される場合、論理積は三つの不等式の同時満足と等価である。一般的には、制約は異なった方法でインスタンス化されることがある。その場合、$\text{共鳴}(t) := C_M \land C_U \land C_V$ は基準として留まり、閾値だけではその定義を尽くさない。

\begin{definition}[共鳴]
\textit{マルチエージェントシステムが時刻 $t$ で共鳴状態にあるとは、意味論的分散条件（$C_M$）、不確実性整合条件（$C_U$）、および価値方向条件（$C_V$）が同時に満たされる場合である。共鳴は差異が複数の次元に渡って首尾一貫して制約される状態である。}
\end{definition}

ここで重要なのは、整合がここで完全な恒等を意味しないことである。$C_M$ は意味的差異が消滅されるのではなく、分散と閾値の観点から制御されることを要求する。同様に、$C_U$ は共有された不確実性の度合いを指し、$C_V$ は評価的標準または判断方向の整合を指す。

従って共鳴は差異が持続する一方で複数次元に渡って首尾一貫して制約される状態である。そのような状態では、意味的差異は混沌的に拡散しないが、特定の構造に従って分布する。この構造は意味状態空間の領域として表現されることがある。
\[
M_i(t) \in B \subseteq \mathbb{R}^d.
\]

ここで $B$ は符号化空間 $\mathbb{R}^d$ の有界領域であり、共鳴状態にあるエージェントの符号化された意味値はこの領域内に分布すると見なされる。この領域は固定点ではない。それは意味が動力学的に変動する範囲を指定する。

動力学系の視点から、そのような領域はアトラクタに類似している。しかし本研究のアトラクタは収束の点ではなく、非収束軌道が制約される領域である。この意味において、共鳴は\textit{非収束アトラクタ}として解釈されることがある。

意味的差異の長期的振る舞いを考慮して、消滅しない長期振幅が残る。規範を最初に取り、次を定義する。
\[
G^{(\infty)}_{ij} := \limsup_{t \to \infty} \|\Delta M_{ij}(t)\|.
\]

非負スカラ写像 $h_{ij}(t) = \|\Delta M_{ij}(t)\|$ に対して、上限の一般的公式は以下を与える。
\[
G^{(\infty)}_{ij} = \lim_{T \to \infty} G_{ij}(T) = \inf_{T \geq 0} \sup_{t \geq T} \|\Delta M_{ij}(t)\|.
\]

有限時刻 $t$ 付きの $\limsup$ を混合する早期ドラフト問題を避けるため、振幅は常に $\|\cdot\|$ と $h_{ij}$ を通じて記述される。

この残存は過去の相互作用によって生成された差異が完全には消滅せず、システム全体に影響し続けることを示す。共鳴状態では、この残存も一定範囲内で制約されている。共鳴は従って現在の状態だけではなく、その歴史を含む全体的構造の安定性を指す。

従って共鳴は非収束下での安定性の条件として位置付けられる。意味は収束することなく構造的に安定のままであることがある。この安定性は制約 $C_M, C_U, C_V$ の同時作用を通じて生じる。

共鳴はセクション3で導入されたコンセンサスの概念を一般化する。コンセンサスは主として意味的差異の大きさと安定性に焦点を当てるのに対し、共鳴は不確実性と価値の次元を追加的に含む。従ってコンセンサスは共鳴の特殊ケースとして理解されることができる。

\subsection{非収束安定性の理論的性質}

本小セクションは、セクション2～4が依拠する数学的基礎を簡潔に列挙する。完全な証明はここでは発展させない。厳密な発展は将来の研究に残される。「定理」という用語は避けられ、さらに仮定を要する主張は明示的に命題または現象論的原理として標識される。特に、更新項 $f$ の有界性だけではそれ自体は軌道の有界性を暗示しない。本小セクションの中心的仮定は従って正不変領域の存在である。

\begin{remark}[有界 $f$ だけは有界軌道を暗示しない]
更新関数 $f$ が一様有界であっても、
\[
M_i(t) = M_i(0) + \int_0^t \sum_j f(\alpha_{ji}(\tau), \Delta M_{ji}(\tau)) d\tau
\]
は一般的には有界である必要がない。$j$ 上の和が長時間にわたって時間依存的な累積で作用する場合、線形成長が生じることがある。従ってそれは有界振幅の安定領域を主張する場合、散逸、射影、または正不変性といったベクトル場についての追加的仮定が不可欠である。

$f$ の一様有界性、
\[
\|f(\alpha, x)\| \leq C \quad \text{for all } \alpha \in [0,1], x \in \mathbb{R}^d,
\]
はただ局所更新が任意に大きくならないことを意味するために取られている。
\end{remark}

\begin{assumption}[正不変有界閉領域]
時間的ベクトル場に対して有界閉集合 $B \subseteq \mathbb{R}^d$ が存在して、容認可能な初期条件から開始する全ての軌道が、一度 $B$ に入れば、永遠に $B$ に留まる。

幾何学的には、境界 $\partial B$ に内向的指示条件を課すことができ、例えば
\[
\left\langle \sum_j f(\alpha_{ji}, \Delta M_{ji}), n_i \right\rangle \leq 0
\]
外向き法線 $n_i$ に対して。本論文は上記の集合水準仮定を採用する。
\end{assumption}

\begin{proposition}[非収束だがかつ振幅有界なレジームの存在]
仮定1が成立し、非収束性公理
\[
\limsup_{t \to \infty} \|\Delta M_{ij}(t)\| > 0
\]
が効力を持つ場合、意味状態は有界領域内に制約されながら厳密な恒等に到達しないことがある。すなわち、有界領域 $B$ と時刻 $T$ が存在して
\[
M_i(t) \in B \quad (\forall t \geq T), \quad \limsup_{t \to \infty} \|\Delta M_{ij}(t)\| > 0.
\]

$f, g, \alpha$ のクラスに対する具体的な十分条件は将来の研究に残される。
\end{proposition}

\textbf{解釈。} 意味が単一のポイントに収束しなくても、正不変領域は $\Delta M$ が内部で非零に留まりながら振幅を制約することができる。二つの層は共存できる。

\textbf{スケッチ。} 正不変性は $\|M_i\|$ の爆発を除外する。一方、非収束性公理は $\|\Delta M\| \to 0$ である軌道を除外する。典型的ケースは単調に収束しない閉または引力的運動を含む。

\begin{assumption}[有限次数]
相互作用グラフ $G(t)$ は有限次数を持つ。すなわち、与えられたエージェントに同時に影響するエージェント数は上から有界である。
\[
\deg(i,t) \leq D \quad \text{for all } i, t.
\]
\end{assumption}

\begin{phenomenological}[疑似収束からの乖離]
非収束システムでは、全ての単一のポイント $M^*$ が正真正銘の勾配収束の極限として機能するわけではない。小さな差異と擾乱は時間とともに累積され、状態は候補ポイント付近に留まるかもしれないが、最終的にそれから乖離する。これは証明された定理ではなく、行動的テンプレートとして提示されている。

\textbf{解釈。} 意味が静的に見えても、公理 $\limsup \|\Delta M\| > 0$ の下では、それは完全な恒等に沈むことはできず、運動に戻ることがある。

\textbf{スケッチ。} 非零差異からの漂流または内部動力学と外部制約を比較する。
\end{phenomenological}

\begin{assumption}[擾乱後の拡張領域]
外部決定論的ノイズ $\eta(t)$ または有界増分を持つ確率過程が状態に単調に作用する擾乱されたシステムに対して、システムは有界領域 $B' \subseteq \mathbb{R}^d$ で正不変に留まると仮定する。ここで $B'$ は元々の $B$ より大きい。これは入力と状態制約に問題を削減する。
\end{assumption}

\begin{phenomenological}[外部擾乱の吸収]
仮定3の下で、擾乱は異常として現れる可能性は少なく、むしろ既存の変動の上に層状化される。軌道はグローバルに $B'$ で制約され、共鳴としての構造は閾値超過によってのみ評価される。
\end{phenomenological}

\begin{proposition}[共鳴状態の記述的意味]
意味的分散 $\text{Var}_E(\Delta M)$ および不確実性と価値発散指示器 $D_U, D_V$ がそれぞれそれらの閾値を下回る場合、$\text{共鳴}(t)$ が満たされる。これはセクション4の定義の直接的な言い換えである。その存在は具体的なモデル係数の割り当てにおいてのみ非自明となる。
\end{proposition}

\subsection{既存理論との相違}

Resonanceverse は単なる「非収束有界軌道」の言い直しや「意味は固定されていない」という主張ではない。その識別的特徴は差異の持続から開始し、そうした差異がいかに意味、不確実性、および価値の複数の条件下で制約され、制度的に扱う可能な構造を形成することができるかを問うことである。

固定点収束モデルでは、安定性はポイントに接近する状態として理解される。そのような枠組みの下では、差異は最終的に削減されるか消滅されるべき何かである。対照的に、本論文は完全な恒等性を前提としない。むしろそれは差異の持続それ自体を構造的条件として扱う。

古典的なコンセンサス形成理論では、中心的概念は一致、平均、またはエージェント状態間の近接である。差異はしばしばコンセンサスプロセス中に減少するべき何かとして扱われる。本論文ではコンセンサスを恒等性ではなく、$\Delta M$ 分布の有界で時間的に安定した形への修正として定義する。

有界軌道とアトラクタについての理論はその軌道の領域内への制約に対処する。しかし、それらは通常、その有界性によって保持されている意味的差異、価値的差異、または不確実性的差異が何かを問わない。また、どの差異が制度的に制約されるかを問わない。Resonanceverse の識別性はそれは有界性そのものにあるのではなく、有界性にその意味を与えるものにある。

カオスおよび複雑性理論は有界だが予測困難な振る舞いに対処する。そのような枠組みでは、差異はシステムの非線形振る舞いとして現れる。この研究が対処するのは、単なる複雑性ではなく、制約可能な $\Delta M$ 構造であり、意味、不確実性、および価値条件を同時に満たしている。

ポスト構造主義の意味理論は非固定性、\textit{différance}、および誤配に対する強力な洞察を提供する。本研究はそうした洞察を否定しない。むしろそれを $\Delta M$ ネットワークおよび共鳴条件として形式化する。従って Resonanceverse の非収束安定性は「有界非収束」ではなく、\textit{条件付き非収束}として位置付けられる。言い換えれば、この理論は単なる有界軌道ではなく、意味的差異、不確実性的差異、および価値的差異が制度的に修正される状態に対処する。

\section{議論：ゲーデル類推と理論的含意}

本セクションはセクション2～4で形式化された意味動力学および非収束安定性を、より広い理論的文脈内に位置付ける。特に本研究の中心的概念——非収束のままで安定する構造——は形式体系における決定不可能性と関連させて再解釈される。この対応はそれは形式的決定不可能性と意味動力学的フロー間の同型ではない。それは厳密に哲学的類推として読まれるべきである。

\subsection{方法論的警告}

\textbf{方法論的警告。} ゲーデル対応は技術的同型ではない。ゲーデルは明示的に決定不可能文を構成し、形式体系からのそれらの独立性を証明する。Resonanceverse は非収束安定性を公理として仮定し、その結果を導出する。類推は形式内部の不完全性と外部的一貫性の\textit{構造的分離}にあり、正式な論理的等価性にはない。証明論的決定不可能性とマルチエージェント $\Delta M$ ネットワークの振幅と分散上の制約は異なる論理的領域に属する。

\subsection{ゲーデル類推}

この視点から、Kurt Gödel の不完全性定理との類推は示唆的である。ゲーデルは、十分に強い形式体系では、体系内で証明可能だが真である命題が存在することを示した。この結果は形式体系の完全性と決定可能性の限界を明らかにした。

ゲーデルの不完全性定理は形式体系内での証明可能性と、その体系外から見た真理評価との分離を可能にする方法を提供する。ここでは、その哲学的含意は意味動力学における収束しないが構造を保持する状態を読む類推として用いられる。

この構造は意味動力学における非収束安定性に対応する。意味的差異 $\Delta M_{ij}(t)$ は時間上で消滅しず、エージェント間の意味は完全に一致することはない。この意味において、意味は非収束的であり、決定不可能である。同時に、しかし、セクション4に示されたように、意味的差異は有界領域内に制約され、共鳴条件下で安定した構造を形成する。

対応は以下のようにまとめられることがある。

\begin{table}[h]
\centering
\begin{tabular}{ll}
\hline
\textbf{意味動力学} & \textbf{形式体系} \\
\hline
意味の非収束性 & 形式体系における決定不可能性 \\
共鳴下の構造安定性 & 証明可能性を超えて存在する真理 \\
$\Delta M$ ネットワーク保存 & 形式体系外の秩序 \\
\hline
\end{tabular}
\end{table}

この視点から、本論文の枠組みは形式体系における証明可能性と真理の分離を、意味動力学における収束と構造安定性の分離として再解釈する。

さらに、この対応は真理の動的再定義を示唆する。現在の枠組みにおいて、意味は固定点として存在しない。それは時間の中で継続的に変動する。共鳴として現れる構造は固定点ではなく、有界領域内で持続するパターンである。この意味において、真理は静的な終点ではなく、動的構造として理解される。

\textit{真理は非収束的な $\Delta M$ 動力学内で持続する安定した構造として解釈されることがある。}

この定式化は真理を到達すべき固定点ではなく、意味の非収束的動力学内で持続する構造として再定義する。それはそれは従来の静的意味論や形式主義的真理概念とは異なる関係的で時間的な理解を提供する。

この再定義はマルチエージェント環境においてとりわけ重要である。ゲーデルの結果が単一の形式体系内の限界に関わるのに対し、現在の枠組みは複数のエージェント間の相互作用を通じて生じる意味的構造に関わる。ここで真理に対応するのは単一の命題ではなく、$\Delta M$ ネットワーク内で安定化した構造である。真理は個々の表現ではなく、関係内に持続するパターンである。

セクション4に示されたように、非収束安定性はまた外部擾乱に対して識別的な堅牢性を持つ。収束システムでは、状態がその固定点で固定されるため、外部ノイズはその状態からの乖離として作用し、全体的構造を不安定化させることがある。収束ポイントが局所疑似解である場合、システムはそこで固定されることがあり、そこから離脱することが困難になる。

対照的に、非収束安定システムでは、意味状態は既に有界領域内で変動しているが。従って外部ノイズは例外的擾乱ではなく、内部変動の一部として吸収される。システムは外部擾乱に対してロバストであり、誤った状態での疑似収束への固定を回避する。このプロパティは意味的構造が誤った状態に安定化することを妨げるメカニズムとして機能する。

これらのプロパティは意味を超えた含意を持ち、制度設計と人工知能へ拡張する。セクション3に示されたように、制度は $\Delta M$ の分布上の制約として理解されることができる。従って制度設計は意味的差異を消滅させるのではなく、それらの分布が制御され、非収束安定構造が維持される条件の設計の問題として再定義される。

同様に、人工知能システムではでは、完全な収束または固定された最適解を目指す設計は、誤った状態での固定と外部擾乱への脆弱性を招くことがある。対照的に、非収束安定性に基づく設計は変動を許可しながら構造を維持することによってより堅牢なシステムを達成することができる。

従って、この枠組みは意味、真理、コンセンサス、制度、およびインテリジェンスを収束的エンティティではなく、非収束的構造として理解するための基礎を提供する。それはポスト構造主義の記述的洞察を保存しながら、それを動的で形式的な理論へ拡張する。結果は、意味を固定的なオブジェクトとして把握するのではなく、関係的で時間的な構造として把握する理論的基礎である。この基礎は数学的精密化、実験的検証、および応用設計を通じた更なる発展の可能性を持つ。

\section{数値検証}

理論的枠組みを検証するため、包括的なマルチエージェントシミュレーションシステムを実装し、セクション2～4およびセクション D に対応する四つのフェーズの実験を実施した。完全な実装は \url{https://github.com/zyx-corporation/resonanceverse_sim} で利用可能である。

\subsection{Phase 1: 二エージェント基礎検証}

\textbf{目的：} 補題1（スケール分離）、セクション2.4おもちゃ例、および公理的基礎を検証する。

\textbf{セットアップ：} セクション2.4で指定されたパラメータを持つ $\mathbb{R}^2$ の二エージェントシステム。シミュレーション時刻：$T_{\max} = 100$ ステップ、$\Delta t = 0.01$。

\textbf{結果：}
\begin{itemize}
\item \textbf{非収束公理：} $\limsup_{t \to \infty} \|\Delta M_{12}(t)\| = 0.1009 > 0$ ✓
\item \textbf{有界性：} $\max_t \|\Delta M_{12}(t)\| = 0.1416 < \varepsilon_1 = 0.15$ ✓
\item \textbf{緩やかな変化：} $\max_t |d\|\Delta M_{12}\|/dt| = 0.0543 < \delta = 0.05$ ✓
\end{itemize}

\textbf{結論：} Phase 1 は意味的差異が有界に留まる一方ゼロに収束しないことを確認し、補題1 を数値的に検証する。補足図1を参照。

\subsection{Phase 2: マルチエージェント・コンセンサス形成}

\textbf{目的：} セクション3コンセンサス定義および制度的制約効果を検証する。

\textbf{セットアップ：} $n = 10$ エージェント、小世界ネットワーク・トポロジー（$k = 4$、$p_{\text{rewire}} = 0.1$）。三つのシナリオ：（A）制約なし、（B）制度的制約（$\text{Var}_E < \theta_M = 0.2$）、（C）$t = 50$ での外部擾乱。

\textbf{結果：}
\begin{itemize}
\item \textbf{シナリオ A（無制約）：} コンセンサス達成 $t = 385$。
\item \textbf{シナリオ B（制約）：} コンセンサス達成 $t = 323$（42\% 高速化）。
\item \textbf{シナリオ C（擾乱）：} システムは 150 ステップ内で共鳴状態に復帰。
\end{itemize}

\textbf{結論：} 制度的制約は $\Delta M$ ネットワーク分布を効果的に規制し、セクション3の制度設計枠組みを確認する。補足図2を参照。

\subsection{Phase 3: 共鳴条件検証}

\textbf{目的：} 定義1および命題1を検証する。

\textbf{セットアップ：} $n = 10$ エージェント、拡張状態空間 $\mathbb{R}^3$（意味 + 不確実性 + 価値）、$T_{\max} = 200$ ステップ。共鳴閾値：$\theta_M = 0.2$、$\theta_U = 0.15$、$\theta_V = 0.25$。

\textbf{結果：}
\begin{itemize}
\item \textbf{共鳴達成：} タイムステップの 98.6\% が $C_M \land C_U \land C_V$ を満たす。
\item \textbf{リアプノフ下降：} $V(t) = \frac{1}{2}\sum_i \|M_i\|^2$ は 58.9\% 減少。
\item \textbf{条件独立性（動力学的）：} $C_M$ と $C_V$ 間の相関：0.609。
\item \textbf{条件独立性（統計的、Phase 4）：} 相関 0.022。
\end{itemize}

\textbf{結論：} 共鳴条件は同時満足可能であり、安定した非収束構造を生成し、定義1および命題1～4を検証する。補足図3を参照。

\subsection{Phase 4: 反証可能性プロトコル検証}

\textbf{目的：} セクション D 反証可能性条件を直接テストで検証する。

\textbf{プロトコル1（非収束公理）：}
\begin{itemize}
\item 複数のドメイン（法的、科学的、創造的）にわたり 100+ 試行を実施。
\item 収束率を測定：$\|\Delta M(t_{\max})\| < 10^{-6}$ （$t_{\max} = 500$ ステップ後）。
\item \textbf{結果：} 300 全試行にわたり 0\% 収束。非収束公理は強く支持される。
\end{itemize}

\textbf{プロトコル3（条件独立性）：}
\begin{itemize}
\item 1000 ランダムエージェント状態を生成。
\item ペアワイズ相関を測定：$\text{corr}(C_M, C_U)$、$\text{corr}(C_M, C_V)$、$\text{corr}(C_U, C_V)$。
\item 独立性棄却閾値：$\max(\text{corr}) > 0.7$。
\item \textbf{結果：} $\max(\text{corr}) = 0.022 \ll 0.7$。条件独立性が確認される。
\end{itemize}

\textbf{結論：} 反証可能性プロトコルは Resonanceverse の科学的性質を確認する。非収束公理は実験的に支持され、共鳴条件の論理的独立性が検証される。補足図4を参照。

\section{結論}

本研究は Resonanceverse を提示した。これは意味を静的な表現ではなく、エージェント間の相互作用を通じて生成される動的構造として再定義する理論的枠組みである。本枠組みでは、意味は固定的なオブジェクトではなく、時間の中で継続的に変動する状態量であり。その変換はエージェント間の差異、すなわち $\Delta M$ によって記述される。

セクション2は意味状態空間、意味的差異 $\Delta M_{ij}(t) := M_j(t) - M_i(t)$、および関係を含む公理系を導入した。それは $\Delta M$ が誤配そのものではなく、誤配プロセスから生じる状態的差異であることを明確にした。本セクションはまた非収束公理を実用的 $\varepsilon$ スケール・コンセンサスおよび共鳴からスケール分離を通じて区別し、厳密な恒等が原理的には達成されないことを明確にした。

セクション3 はフレームワークをマルチエージェント・システムに拡張し、意味論的構造を意味的差異のネットワークとして記述した。それはコンセンサスを恒等性ではなく、意味的差異の分布が有界で時間的に安定する状態として再定義した。従ってコンセンサスは差異の消滅ではなく、差異の構造的修正として示された。

セクション4 は非収束安定性を非収束状態の安定性概念として導入し、共鳴を $C_M, C_U, C_V$ の同時満足として定義した。それは意味が収束することなく安定した構造を形成することがあり得ることを明らかにする理論的足場を整理した。振幅有界性と非収束性が仮定1などの条件下で共存できることを示した。

セクション5 はフレームワークをより広い理論的文脈に位置付け、形式体系における決定不可能性に特に関連させた。意味動力学における非収束と構造安定性の分離は証明可能性と真理の分離に類推として解釈された。これは真理が固定点ではなく、非収束的動力学内の安定した構造として再解釈される視点を導入した。

セクション6 は四つのフェーズにわたる包括的な数値検証を提示した。Phase 1～4 実験は全ての理論的予測を確認する。非収束公理（300 試行にわたり 0\% 収束）、共鳴条件満足（タイムステップの 98.6\%）、リアプノフ安定性（58.9\% 下降）、および共鳴条件の論理的独立性（$\max(\text{corr}) = 0.022$）。

\subsection{意味の再定義}

これらの結果から、本論文は意味の以下の再定義を提供する。

\textit{意味は主体に内在する固定的内容ではない。それはエージェント間の差異のネットワークとして存在し、そのような差異の時間的累積を通じて形成される。その安定性は収束によってではなく、差異を保持しながら制約された範囲内で持続することによって成立する。}

この再定義はセマンティクスだけでなく、コンセンサス形成、制度設計、および人工知能にも含意を持つ。特に、収束志向モデルを非収束安定性で置き換えることは、外部擾乱に対する堅牢性と誤った固定の回避に対して有用な枠組みを提供することができる。

\subsection{反証可能性と限定条件}

本研究は意味が原理的に完全な恒等性に収束しないという非収束公理に根ざしている。しかし、この公理は全ての対象に無条件に適用される実証的事実として主張されない。従って、この理論は以下の条件下で改訂または棄却されることがある。

\textbf{条件1：} 複数の意味的ドメイン全体にわたって、符号化された意味的差異 $\Delta M_{ij}(t)$ が十分に長い相互作用の後、安定的にゼロに収束することが反復可能に示される場合、非収束公理を一般原理として維持することはできず、そのスコープを限定する必要がある。

\textbf{条件2：} $\Delta M$ の分布が、制度的制約または文脈的親和性係数 $\alpha_{ij}(t)$ なしに、一般的に有界で安定することが示される場合、$\Delta M$ 分布への制約として制度を位置付ける必要性は弱まる。

\textbf{条件3：} 意味、不確実性、および価値の三つの条件——$C_M, C_U, C_V$——が独立として分離されず、代わりに単一のスカラ指標に削減されることができる場合、その論理的論理積として共鳴を定義する有意性は減少する。

\textbf{条件4：} $\Delta M$ が観測可能またはも推定可能な量として安定的に構成されない場合、理論は形式モデルとしての堅牢性を失う。

\textbf{条件5：} 共鳴条件を満たす状態が、実際のマルチエージェント・システムで安定性ではなく抑圧、均質化、または病的閉鎖を生成する場合、共鳴は望ましい安定性ではなく疑似安定性として再定義されなければならない。

これらの限定条件は理論の単なる弱点ではなく、将来の研究課題でもある。Resonanceverse は完成した閉じられたシステムではない。それは $\Delta M$ の観測、制度的制約の設計、および共鳴条件の定量化を通じて検査されなければならない開放的な理論的枠組みである。

\subsection{将来の仕事}

将来の仕事はここで提示された枠組みの数学的精密化を含む。具体的には、$\Delta M$ を確率過程として形式化し、共鳴条件を定量的に評価し、非収束安定性の存在と一意性に関する定理を導出することが必要であろう。追加的なマルチエージェント・シミュレーションおよび実際の意味的ドメインにおける実験的検証を通じて理論をテストすることもまた重要である。

制度および人工知能システムへの理論応用では、中心的設計問題は $\Delta M$ の分布をいかに制御するか、およびどの条件の下で共鳴が維持されるかであろう。その点で本研究は意味動力学に根ざした新たな設計原則を提供する。

本論文は意味を収束的オブジェクトではなく、非収束的でかつ構造的に安定したプロセスとして理解するための理論的基礎を提示した。この視点は意味の再構築だけでなく、複雑なマルチエージェント・システムにおける安定性の理解へも新しい枠組みを提供する。

簡潔に、本研究の立場は以下のように述べられることがある。

\textit{意味は恒等性を通じて生じない。それは差異を保持しながら持続する構造として現れる。この原則は セマンティクスの再構築のためだけでなく、複雑なマルチエージェント・システムにおける安定性の理解のためのも指針を提供する。}

\begin{thebibliography}{99}

\bibitem{azuma1998} 青山 浩（1998）「存在論的、郵便的——ジャック・デリダについて」新潮社。

\bibitem{berger1966} P. L. Berger and T. Luckmann, \textit{The Social Construction of Reality: A Treatise in the Sociology of Knowledge}, Doubleday, 1966.

\bibitem{degroot1974} M. H. DeGroot, ``Reaching a consensus,'' \textit{Journal of the American Statistical Association}, vol. 69, no. 345, pp. 118--121, 1974.

\bibitem{derrida1967} J. Derrida, \textit{De la Grammatologie}, Les Éditions de Minuit, 1967.

\bibitem{derrida1972} J. Derrida, ``La différance,'' in \textit{Marges de la Philosophie}, pp. 1--29, Les Éditions de Minuit, 1972.

\bibitem{godel1931} K. Gödel, ``Über formal unentscheidbare Sätze der Principia Mathematica und verwandter Systeme I,'' \textit{Monatshefte für Mathematik und Physik}, vol. 38, pp. 173--198, 1931.

\bibitem{khalil2002} H. K. Khalil, \textit{Nonlinear Systems} (3rd ed.), Prentice Hall, 2002.

\bibitem{landauer1997} T. K. Landauer and S. T. Dumais, ``A solution to Plato's problem: The latent semantic analysis theory of acquisition, induction, and representation of knowledge,'' \textit{Psychological Review}, vol. 104, no. 2, pp. 211--240, 1997.

\bibitem{mikolov2013} T. Mikolov, K. Chen, G. Corrado, and J. Dean, ``Efficient estimation of word representations in vector space,'' arXiv:1301.3781, 2013.

\bibitem{moreau2005} L. Moreau, ``Stability of multiagent systems with time-dependent communication links,'' \textit{IEEE Transactions on Automatic Control}, vol. 50, no. 2, pp. 169--182, 2005.

\bibitem{north1990} D. C. North, \textit{Institutions, Institutional Change and Economic Performance}, Cambridge University Press, 1990.

\bibitem{olfati2007} R. Olfati-Saber, J. A. Fax, and R. M. Murray, ``Consensus and cooperation in networked multi-agent systems,'' \textit{Proceedings of the IEEE}, vol. 95, no. 1, pp. 215--233, 2007.

\bibitem{shannon1948} C. E. Shannon, ``A mathematical theory of communication,'' \textit{The Bell System Technical Journal}, vol. 27, no. 3, pp. 379--423; vol. 27, no. 4, pp. 623--656, 1948.

\bibitem{sontag1989} E. D. Sontag, ``Smooth stabilization implies coprime factorization,'' \textit{IEEE Transactions on Automatic Control}, vol. 34, no. 4, pp. 435--443, 1989.

\bibitem{strogatz2015} S. H. Strogatz, \textit{Nonlinear Dynamics and Chaos: With Applications to Physics, Biology, Chemistry, and Engineering} (2nd ed.), Westview Press, 2015.

\end{thebibliography}

\end{document}
